\begin{tabularx}{\textwidth}{L{33mm}X X}
\toprule
\textbf{Probléma} & \textbf{Lényege} & \textbf{Következmény / kezelés} \\
\midrule
Egyértelműség hiánya & A természetes nyelv többértelmű, ugyanazt a mondatot az ügyfél és a fejlesztő eltérően értelmezheti. & Pontosítás, példák, strukturált természetes nyelv, modellek, prototípusok és validáció szükséges. \\
Követelmények keveredése & Funkcionális, nemfunkcionális, szakterületi, üzleti célok és tervezési részletek összecsúszhatnak. & A követelményeket csoportosítani, címkézni és külön szinteken dokumentálni kell. \\
Követelmények ötvöződése & Több eltérő igény egyetlen hosszú követelményként jelenik meg. & Szét kell bontani ellenőrizhető, priorizálható, tesztelhető követelményekre. \\
Teljesség és konzisztencia hiánya & Nagy rendszereknél könnyű kihagyni szolgáltatásokat, illetve ellentmondó elvárásokat rögzíteni. & Követelményellenőrzés, áttekintés, érintetti egyeztetés és változáskövetés szükséges. \\
Nehezen verifikálható elvárások & Különösen a nemfunkcionális követelmények gyakran homályosak, például „legyen könnyen használható”. & Objektív metrikákkal kell megfogalmazni: például válaszidő, hibaarány, rendelkezésre állás, megtanulási idő. \\
Szakterületi félreértések & A szakértők természetesnek vesznek olyan ismereteket, amelyek a fejlesztőknek nem egyértelműek. & Szójegyzék, szakterületi modell, interjú, megfigyelés és forgatókönyv segíthet. \\
\bottomrule
\end{tabularx}
